% !TEX root = ../Monografia.tex
% ---------------------------------------------------------------------------------------------------------------------------------------------------------
\section{Procedimentos de Teste e Coleta de Dados}
% ---------------------------------------------------------------------------------------------------------------------------------------------------------

Esta seção descreve os procedimentos utilizados para validar funcionalmente o protótipo \texttt{collector\_flutter}. O objetivo principal dos testes foi observar três aspectos fundamentais do sistema: (i) a coerência das métricas coletadas, (ii) o impacto percebido da coleta durante a execução e (iii) a pertinência das recomendações heurísticas geradas pelo sistema.

Para garantir validade científica e reprodutibilidade, os testes foram conduzidos em ambiente controlado, seguindo protocolos experimentais inspirados nas diretrizes de experimentação em engenharia de software propostas por Wohlin \textit{et al.}~\cite{wohlin2012}.

\subsection{Protocolo de Execução dos Testes}

Os experimentos foram executados em modo \textit{profile} do Flutter, pois esse modo permite medições mais próximas do comportamento real da aplicação, eliminando interferências associadas ao modo \textit{debug}. Antes de cada execução experimental, o dispositivo foi reiniciado e o aplicativo executado em estado limpo, garantindo consistência entre as medições.

O protocolo de execução foi composto pelas seguintes etapas:

\begin{enumerate}

\item \textbf{Preparação do ambiente}: reinicialização do dispositivo e encerramento de processos em segundo plano, assegurando estabilidade do ambiente de execução.

\item \textbf{Inicialização do coletor}: ativação do módulo \texttt{ResourceCollector} por meio do método \texttt{start()}, responsável por iniciar a captura das métricas de desempenho.

\item \textbf{Execução do cenário experimental}: cada cenário foi acionado no aplicativo de exemplo por tempo suficiente para observar a reação do coletor e do painel visual.

\item \textbf{Registro das métricas}: durante a execução do cenário, o coletor registrou continuamente métricas de desempenho relacionadas à renderização, uso de memória, chamadas de rede e reconstruções de widgets.

\item \textbf{Encerramento da coleta}: ao término do experimento, o método \texttt{stop()} foi acionado e os dados coletados foram exportados para arquivos estruturados no formato JSON.

\end{enumerate}

Os cenários foram executados manualmente em condições controladas, com o objetivo de verificar a operacionalidade do sistema e identificar inconsistências nas heurísticas. Os resultados apresentados neste trabalho são de natureza qualitativa.

\subsection{Definição dos Cenários Experimentais}

Quatro cenários experimentais foram definidos com o objetivo de simular diferentes tipos de carga comuns em aplicações móveis. Cada cenário foi projetado para provocar um tipo específico de comportamento de desempenho.

\begin{itemize}

\item \textbf{Cenário 1 – Reconstruções intensivas de interface}

Neste cenário, um conjunto de widgets foi reconstruído repetidamente em alta frequência, simulando aplicações com atualização constante de interface. Esse cenário foi utilizado para avaliar o impacto das reconstruções de widgets sobre o tempo de renderização e o FPS.

\item \textbf{Cenário 2 – Chamadas de rede simultâneas}

Foram disparadas múltiplas requisições HTTP paralelas, utilizando o método \texttt{recordNetwork()}, com o objetivo de observar o comportamento do sistema sob carga de rede e verificar a capacidade do coletor de registrar eventos relacionados a comunicação externa.

\item \textbf{Cenário 3 – Alocação intensiva de memória}

Objetos de grande volume foram criados e descartados continuamente durante a execução do teste. Esse cenário permitiu analisar padrões de uso de memória e observar eventos de coleta de lixo (\textit{garbage collection}).

\item \textbf{Cenário 4 – Carga combinada}

Neste cenário, os três tipos de carga anteriores foram executados simultaneamente, simulando um comportamento mais próximo de aplicações reais, nas quais operações de interface, rede e memória ocorrem de forma concorrente.

\end{itemize}

Esses cenários permitiram avaliar o comportamento do sistema em diferentes condições de uso, identificando possíveis limitações e avaliando a robustez da ferramenta.

\subsection{Coleta de Métricas}

As métricas foram coletadas automaticamente pelo módulo \texttt{Collector}, implementado inteiramente em Dart. A captura dos dados foi realizada em tempo real utilizando \texttt{Streams} internas do Flutter.

As métricas analisadas foram:

\begin{itemize}

\item \textbf{FPS (Frames por Segundo)}: indicador da fluidez da renderização da interface;

\item \textbf{Tempo médio de renderização de quadro}: calculado a partir dos eventos de \texttt{FrameTiming};

\item \textbf{Número de reconstruções de widgets}: contabilizado por meio do componente \texttt{RebuildObserver};

\item \textbf{Uso de memória}: obtido via API \texttt{vmService.getAllocationProfile()};

\item \textbf{Número de chamadas de rede}: monitorado por meio da instrumentação do método \texttt{recordNetwork()}.

\end{itemize}

Os dados coletados foram serializados em formato JSON estruturado, permitindo inspeção posterior e reprodutibilidade dos registros observados.

\subsection{Métodos de Comparação}

Para apoiar a verificação das medições obtidas pelo protótipo, os resultados foram confrontados qualitativamente com ferramentas oficiais de análise de desempenho do ecossistema Android.

As ferramentas utilizadas como referência foram:

\begin{itemize}

\item \textbf{Android Profiler}, utilizado para medir consumo de CPU, memória e FPS;

\item \textbf{ADB (\texttt{adb shell dumpsys gfxinfo})}, utilizado para obter métricas detalhadas de renderização de quadros.

\end{itemize}

A comparação quantitativa por erro percentual não foi executada nesta versão, mas a expressão abaixo documenta o cálculo previsto para estudos futuros:

\[
Erro (\%) = \frac{|M_{ref} - M_{proto}|}{M_{ref}} \times 100
\]

onde \(M_{ref}\) representa o valor obtido pela ferramenta de referência e \(M_{proto}\) representa o valor registrado pelo protótipo.

\subsection{Análise dos Resultados}

Os dados coletados foram analisados de forma qualitativa, considerando a coerência entre o cenário acionado, as métricas exibidas no painel e as recomendações emitidas pelo módulo \texttt{Recommender}.

Além disso, o \textit{overhead} foi tratado como critério metodológico para avaliações futuras, comparando o tempo médio de renderização de quadros com o sistema ativado e desativado.

O cálculo previsto utiliza a seguinte expressão:

\[
Overhead (\%) = \frac{T_{com} - T_{sem}}{T_{sem}} \times 100
\]

onde \(T_{com}\) representa o tempo médio de execução com o coletor ativado e \(T_{sem}\) corresponde ao tempo de execução sem monitoramento.

\subsection{Controle Experimental e Reprodutibilidade}

Para garantir a validade experimental dos resultados, foram adotadas medidas de controle que minimizam interferências externas durante os testes:

\begin{itemize}

\item execução em modo \textit{profile};

\item utilização das mesmas versões do Flutter SDK e dependências em todas as execuções;

\item execução dos testes em ambiente com temperatura e carga de bateria estáveis;

\item repetição manual dos cenários quando necessário para confirmar o comportamento observado;

\item padronização do formato de saída das métricas em arquivos JSON.

\end{itemize}

Essas práticas auxiliam a reprodução dos cenários por outros pesquisadores, reforçando a transparência da avaliação funcional do protótipo \texttt{collector\_flutter}.
\begin{table}[H]
\centering
\caption{Resumo dos cenários experimentais utilizados nos testes.}
\begin{tabular}{|c|p{4cm}|p{6cm}|}
\hline
\textbf{Cenário} & \textbf{Tipo de Carga} & \textbf{Objetivo} \\
\hline
1 & Reconstruções de interface & Avaliar impacto de rebuilds frequentes na renderização \\
\hline
2 & Chamadas de rede simultâneas & Observar comportamento sob múltiplas requisições HTTP \\
\hline
3 & Alocação intensiva de memória & Identificar padrões de consumo e eventos de garbage collection \\
\hline
4 & Carga combinada & Simular comportamento próximo de aplicações reais \\
\hline
\end{tabular}
\end{table}
